\documentclass{article}
\usepackage{amsmath, amssymb, amsthm}

\setlength{\parindent}{0pt}

\newtheorem*{claim}{Claim}

\begin{document}

\begin{claim}
    $\gcd(ka, kb) = k \cdot \gcd(a, b)$ for all $k > 0$.
\end{claim}

\begin{proof}
    Let $d = \gcd(a, b)$. Then $d \mid a$ and $d \mid b$, so $kd \mid ka$ and $kd \mid kb$. Hence $kd$ is a common
    divisor of $ka$ and $kb$.

    \medskip

    Since $d = \gcd(a, b)$, there exist integers $x, y$ such that:
    \[
    ax + by = d
    \]
    Multiplying by $k$:
    \[
    kax + kby = kd
    \]
    Let $c$ be an arbitrary common divisor of $ka$ and $kb$. Then $c | kax$ and $c | kby$, so $c | kd$. Therefore 
    every common divisor of $ka$ and $kb$ divides $kd$, and hence is at most $kd$. Since $kd$ itself is a common 
    divisor, it is the greatest common divisor.
    
    \medskip

    Hence $\gcd(ka, kb) = k \cdot \gcd(a, b)$.
\end{proof}

\end{document}
